\documentclass[11pt,a4paper]{article}
\usepackage[margin=27mm]{geometry}
\usepackage{amsmath,amssymb,amsthm,hyperref}
\newtheorem{theorem}{Theorem}
\newcommand{\gen}[1]{\langle#1\rangle}
\newcommand{\Sz}{\operatorname{Sz}}
\title{An algebraic local action does not force thinness}
\author{Human Compression project}\date{}
\begin{document}\maketitle
\begin{abstract}
We give three explicit involutions in $\Sz(8)$ whose parabolics have
orders $14,10,26$. The resulting rank-three coset geometry satisfies
RWPRI and the algebraic definition of local $2$-transitivity used in
the fixed source. Its actual panels have three elements. A short matrix
identity exhibits a third chamber, and a separate literal certificate
checks every panel. This refutes the $q=8$ thinness clause of the source's
Theorem~6.1; it makes no assertion about the separate $q>8$ clause.
\end{abstract}
\section{A small matrix certificate}
Work in $F=\mathbb F_2[X]/(X^3+X+1)$. Integers from zero to seven encode
coefficients in the basis $1,X,X^2$, so $2=X$, $6=X^2+X$, and
addition is exclusive-or. Put
\[
 a=\begin{pmatrix}0&0&0&1\\0&0&1&0\\0&1&0&0\\1&0&0&0\end{pmatrix},\quad
 b=\begin{pmatrix}6&2&0&1\\2&0&1&0\\7&5&0&2\\7&7&2&6\end{pmatrix},\quad
 c=\begin{pmatrix}1&1&0&1\\1&0&1&0\\1&0&0&1\\1&1&1&1\end{pmatrix}.
\]
These are elements of the classical Suzuki matrix group. To make their
membership explicit, use $w=a$, $z(t)=u(0,t)$, and
\[
u(s,t)=\begin{pmatrix}
1&0&0&0\\s&1&0&0\\s^5+t&s^4&1&0\\s^6+st+t^4&t&s&1
\end{pmatrix}.
\]
Then $b=z(2)wz(2)$ and
$c=u(1,1)wz(1)wu(1,1)^{-1}$. The classical group model and the
maximality of $D_{14}$ in $\Sz(8)$ are explicit structural premises.
All the calculations below use ordinary multiplication of these matrices.

The three matrices are distinct nonidentity involutions. Direct multiplication
gives
\[
 |ab|=7,\qquad |ac|=5,\qquad |bc|=13,\qquad
 (ac)^2(ab)^3a=(bc)^6.                                      \tag{1}
\]
For the order statements it is enough to check the indicated prime power
is the identity and the product itself is not. An independent reader
performs these checks by polynomial multiplication modulo $X^3+X+1$.
Thus (1) is a small certificate that can be checked without constructing
the ambient group.

\section{The geometry and the extra chamber}
Let $G=\Sz(8)$ and put
\[
 G_0=\gen{a,b},\qquad G_1=\gen{a,c},\qquad G_2=\gen{b,c}.
\]
We use left cosets; inversion identifies them with the source's right
cosets and preserves incidence and common representatives.
Two involutions whose product has odd prime order $p$ generate $D_{2p}$.
Consequently the orders of our parabolics are $14,10,26$. The gcd of
any two orders is two; since the corresponding parabolics share their
indicated involution, we obtain
\[
 G_0\cap G_1=\gen a,\quad G_0\cap G_2=\gen b,\quad
 G_1\cap G_2=\gen c,\quad G_0\cap G_1\cap G_2=1.           \tag{2}
\]
The last equality follows from $a\ne b$.

\begin{theorem}
The coset geometry $\Gamma(G;G_0,G_1,G_2)$ is RWPRI and locally
$2$-transitive under the literal algebraic-residue convention of L114.
It is not thin and is not flag transitive.
\end{theorem}
\begin{proof}
Two incident cosets have a common representative, which extends the
pair to a chamber by its coset of the missing type. The same observation
extends smaller flags. Hence this rank-three coset pregeometry is a geometry.

For the empty flag, $G_0=D_{14}$ is maximal in $G$. Within each parabolic,
each intersection subgroup in (2) has prime index, respectively seven,
five or thirteen, and hence is maximal. The trivial triple intersection
is maximal in each order-two pair intersection. These exhaust the
subsets required by the source's RWPRI definition. Its algebraic
rank-one residues are precisely the regular actions $C_2/1$ on two
points, so are $2$-transitive.

Consider the base edge $(G_1,G_2)$. The cosets $G_0$ and $cG_0$ give
its two standard extensions. There is also the coset $xG_0$, where
$x=(ac)^2$. It meets $G_1$ at $x$, and (1) shows that it meets $G_2$
at $(bc)^6$. It is distinct from $G_0$: an element of order five cannot
belong to $G_0\cap G_1=\gen a$.

It is distinct from $cG_0$ as well. If they were equal, then
$cx\in G_0\cap G_1=\gen a$. The equation $cx=1$ is impossible
since $x$ has order five and $c$ has order two. The equation $cx=a$
would give $(ac)^2=ca=(ac)^{-1}$, contradicting $|ac|=5$.
There are therefore at least three chambers over this edge, disproving
thinness. The extra chamber has no common representative: such a
representative would lie in $G_1\cap G_2=\{1,c\}$, which would identify
its third coset with one of the two standard cosets. A translate of the
base chamber always has a common representative. Thus the extra chamber
also disproves flag transitivity.
\end{proof}

\section{Complete panels and exact source scope}
The retained finite data contain the full lists of the three small
parabolics. The reader verifies every group product, inverse and generated
element, independently of the acquisition code. For missing type $i$,
every coset incident with $G_j$ has a representative in $G_j$.
Listing the distinct $gG_i$, $g\in G_j$, and retaining those meeting
$G_k$ therefore checks the entire actual panel. In this example all
three panels have exactly three elements, with representatives
\[
\begin{array}{c|c|c}
\text{missing type}&\text{standard representatives}&\text{extra representative}\\\hline
0&1,c&(ac)^2\\
1&1,b&(ab)^3\\
2&1,a&(ba)^3
\end{array}
\]
Each retained coset comes with its full intersections with both base
parabolics and with their intersection. Only the two standard cosets
meet the latter. The short proof above needs only one extra coset;
the complete panel calculation supplies the exact sizes.

The source is D.~Leemans, \emph{The rank 3 geometries of the simple
Suzuki groups $\Sz(q)$}, Note di Matematica 19(1) (1999), 43--63.
Printed p.~45 defines the residue at a subflag of the distinguished
chamber from intersections of parabolics. Printed p.~46 uses that
residue in its local $2$-transitivity condition. It separately defines
thinness by the number of actual chambers over an edge and explicitly
qualifies agreement with earlier definitions by flag transitivity.
The $q=8$ clause of Theorem~6.1 on p.~61 claims that its rank-three
geometries are thin. The theorem above refutes that literal clause.
It is not a counterexample to the different assertion obtained by
requiring flag transitivity or $2$-transitivity on the full actual panels.

The discovery run stopped after the seventy-first prescribed candidate;
its unmodified prefix and terminal receipt are retained. The independent
reader checks this one witness, not the discovery prefix or a classification
of all embeddings. The nine earlier positive markings remain valid.
The rank bound and the three rank-two classes for $q=8$ are independently
accepted elsewhere. The separate thinness clause for prime $q-1>7$
remains an obligation, as do the final source joins.

\paragraph{Attribution and credit.}
The group model, incidence definitions and maximal-subgroup premise are
borrowed as stated. The explicit witness, identity (1), source comparison
and independent finite reader are the contribution of this component.
No priority or later-edition claim is made. This resolves one additional
claim in an original whose principal refutation is already registered;
it adds no original, principal-refutation, conjecture, Lean or compression
count. General source application still requires distinct mathematical review.
\end{document}
